\documentclass[11pt,a4paper]{article}

\usepackage[utf8]{inputenc}
\usepackage[indonesian]{babel}
\usepackage[margin=2cm]{geometry}
\usepackage{amsmath,amssymb,amsfonts}
\usepackage{booktabs}
\usepackage{longtable}
\usepackage{multirow}
\usepackage{enumitem}
\usepackage{titlesec}
\usepackage{float}
\usepackage{xcolor}
\usepackage{tcolorbox}
\usepackage{fancyhdr}
\usepackage{hyperref}

\hypersetup{
    colorlinks=true,
    linkcolor=blue,
    urlcolor=blue,
    pdftitle={Pembahasan Lengkap 20 Soal Latihan Persiapan Diskusi 1},
    pdfauthor={MA4151 Kriptografi}
}

\pagestyle{fancy}
\fancyhf{}
\rhead{\textbf{MA4151 Kriptografi}}
\lhead{Solusi Lengkap 20 Soal Latihan Diskusi 1}
\rfoot{Halaman \thepage}
\setlength{\headheight}{14pt}

% Color Definitions
\definecolor{primary}{RGB}{30, 55, 153}
\definecolor{boxheader}{RGB}{30, 55, 153}
\definecolor{boxbg}{RGB}{248, 249, 250}

\definecolor{colklasik}{RGB}{25, 85, 155}
\definecolor{colbgklasik}{RGB}{242, 247, 253}

\definecolor{colanalisis}{RGB}{140, 40, 110}
\definecolor{colbganalisis}{RGB}{252, 244, 250}

\definecolor{colshannon}{RGB}{38, 130, 70}
\definecolor{colbgshannon}{RGB}{243, 250, 245}

\tcbset{
    arc=1.5mm,
    left=3.5mm,
    right=3.5mm,
    top=3mm,
    bottom=3mm,
    fonttitle=\bfseries
}

% Problem Boxes
\newtcolorbox{soalklasik}[1]{
    colback=colbgklasik,
    colframe=colklasik,
    title={#1}
}

\newtcolorbox{soalanalisis}[1]{
    colback=colbganalisis,
    colframe=colanalisis,
    title={#1}
}

\newtcolorbox{soalshannon}[1]{
    colback=colbgshannon,
    colframe=colshannon,
    title={#1}
}

\newenvironment{solusi}{\noindent\textbf{\textsf{Penyelesaian:}}\par\vspace{0.2em}}{\hfill$\blacksquare$\vspace{1em}\par\hrule\vspace{1em}}

\title{\vspace{-1.2cm}\textbf{\Large SOLUSI \& PEMBAHASAN MATEMATIS LENGKAP}\\\large 20 Soal Latihan Persiapan Diskusi Kelompok 1 (Kriptografi Klasik, Kriptoanalisis, \& Teori Shannon Inti)\\\large\textbf{Mata Kuliah MA4151 Kriptografi}}
\author{\textbf{Disusun dengan Notasi \& Simbol Matematika Formal}}
\date{Semester I -- Menjelang Diskusi Kelompok 1 Oktober}

\begin{document}

\maketitle
\vspace{-0.5cm}

\begin{abstract}
Dokumen ini memuat \textbf{pembahasan matematis lengkap} untuk ke-20 butir soal latihan persiapan Diskusi Kelompok 1 MA4151 Kriptografi. Setiap soal dituliskan kembali secara utuh, diikuti langkah penurunan simbolik formal, persamaan aljabar terstruktur, serta pembuktian deduktif langkah demi langkah sesuai standar penulisan tugas matematika ITB.
\end{abstract}

\vspace{0.2cm}
\hrule
\vspace{0.3cm}

\tableofcontents
\newpage

% =========================================================================
\section{Solusi Bagian I: Kriptografi Klasik (Soal 1--8)}
% =========================================================================

\begin{soalklasik}{Soal 1 [Klasik - Shift]}
Pandang sandi geser atas alfabet $\mathbb{Z}_{26}$ dengan fungsi enkripsi $e_k(x) \equiv x + k \pmod{26}$ dan dekripsi $d_k(y) \equiv y - k \pmod{26}$.
\begin{enumerate}[label=(\alph*)]
    \item Buktikan secara formal bahwa himpunan fungsi enkripsi $\mathcal{E} = \{e_k : k \in \mathbb{Z}_{26}\}$ terhadap operasi komposisi fungsi $(\circ)$ membentuk grup Abelian berorde 26 yang isomorfik terhadap grup aditif $(\mathbb{Z}_{26}, +)$.
    \item Berdasarkan sifat grup pada bagian (a), buktikan bahwa penerapan enkripsi berulang menggunakan dua kunci berturut-turut $k_1, k_2 \in \mathbb{Z}_{26}$ tidak meningkatkan tingkat keamanan sistem kripto.
\end{enumerate}
\end{soalklasik}

\begin{solusi}
\textbf{(a)} Tinjau himpunan $\mathcal{E} = \{e_k : k \in \mathbb{Z}_{26}\}$ dengan operasi komposisi fungsi $(\circ)$:
\begin{enumerate}[leftmargin=*]
    \item \textbf{Sifat Tertutup:} Untuk sebarang $e_{k_1}, e_{k_2} \in \mathcal{E}$ dan untuk setiap $x \in \mathbb{Z}_{26}$:
    \begin{align*}
    (e_{k_2} \circ e_{k_1})(x) &= e_{k_2}(e_{k_1}(x)) \\
    &\equiv (x + k_1) + k_2 \pmod{26} \\
    &\equiv x + (k_1 + k_2 \bmod 26) \pmod{26} \\
    &= e_{(k_1 + k_2 \bmod 26)}(x).
    \end{align*}
    Karena $(k_1 + k_2 \bmod 26) \in \mathbb{Z}_{26}$, maka $(e_{k_2} \circ e_{k_1}) \in \mathcal{E}$.
    
    \item \textbf{Sifat Asosiatif:} Berlaku secara umum untuk komposisi pemetaan fungsi:
    \[
    \forall e_{k_1}, e_{k_2}, e_{k_3} \in \mathcal{E}, \quad (e_{k_3} \circ e_{k_2}) \circ e_{k_1} = e_{k_3} \circ (e_{k_2} \circ e_{k_1}).
    \]
    
    \item \textbf{Elemen Identitas:} Terdapat $e_0 \in \mathcal{E}$ sedemikian hingga $\forall x \in \mathbb{Z}_{26}$:
    \[
    e_0(x) \equiv x + 0 \equiv x \pmod{26}.
    \]
    Maka $\forall e_k \in \mathcal{E}$, $e_k \circ e_0 = e_0 \circ e_k = e_k$.
    
    \item \textbf{Elemen Invers:} Untuk setiap $e_k \in \mathcal{E}$, pilih $e_{26-k} \in \mathcal{E}$ (untuk $k=0$, inversnya $e_0$). Maka:
    \[
    (e_{26-k} \circ e_k)(x) \equiv x + k + (26 - k) \equiv x + 26 \equiv x \pmod{26} = e_0(x).
    \]
    
    \item \textbf{Sifat Komutatif (Abelian):}
    \[
    (e_{k_1} \circ e_{k_2})(x) \equiv x + k_2 + k_1 \equiv x + k_1 + k_2 \equiv (e_{k_2} \circ e_{k_1})(x) \pmod{26}.
    \]
    
    \item \textbf{Isomorfisma $(\mathcal{E}, \circ) \cong (\mathbb{Z}_{26}, +)$:} Definisikan pemetaan $\Phi: (\mathbb{Z}_{26}, +) \to (\mathcal{E}, \circ)$ dengan $\Phi(k) = e_k$.
    \begin{itemize}
        \item Homomorfisma: $\Phi(k_1 + k_2) = e_{(k_1 + k_2 \bmod 26)} = e_{k_2} \circ e_{k_1} = \Phi(k_1) \circ \Phi(k_2)$.
        \item Bijektif: $\ker(\Phi) = \{k \in \mathbb{Z}_{26} : e_k = e_0\} = \{0\}$, sehingga $\Phi$ injektif. Karena $|\mathbb{Z}_{26}| = |\mathcal{E}| = 26$, maka $\Phi$ surjektif.
    \end{itemize}
\end{enumerate}
$\therefore (\mathcal{E}, \circ)$ adalah grup Abelian berorde 26 dan $(\mathcal{E}, \circ) \cong (\mathbb{Z}_{26}, +)$.

\textbf{(b)} Misalkan pesan $p \in \mathbb{Z}_{26}$ dienkripsi dengan kunci $k_1$, lalu dienkripsi kembali dengan kunci $k_2$:
\[
c = (e_{k_2} \circ e_{k_1})(p) = e_{k_3}(p), \quad \text{di mana } k_3 \equiv k_1 + k_2 \pmod{26}.
\]
Karena $(\mathcal{E}, \circ)$ tertutup, pemetaan komposit tersebut ekuivalen dengan enkripsi sandi geser tunggal menggunakan kunci $k_3 \in \mathbb{Z}_{26}$. Ukuran ruang kunci efektif penyerang tetap:
\[
|\mathcal{K}_{\text{efektif}}| = |\mathbb{Z}_{26}| = 26.
\]
$\therefore$ Enkripsi berganda tidak memperluas ruang kunci dan tidak meningkatkan keamanan.
\end{solusi}

\begin{soalklasik}{Soal 2 [Klasik - Affine]}
Diberikan sandi Affine atas $\mathbb{Z}_{26}$ dengan fungsi enkripsi $e_K(x) \equiv ax + b \pmod{26}$ untuk suatu kunci $K = (a, b) \in \mathbb{Z}_{26} \times \mathbb{Z}_{26}$.
\begin{enumerate}[label=(\alph*)]
    \item Buktikan bahwa fungsi enkripsi $e_K$ merupakan pemetaan bijektif jika dan hanya jika $\gcd(a, 26) = 1$.
    \item Turunkan bentuk eksplisit dari fungsi dekripsi $d_K(y)$ dan buktikan secara aljabar bahwa $d_K(e_K(x)) \equiv x \pmod{26}$ untuk setiap $x \in \mathbb{Z}_{26}$.
\end{enumerate}
\end{soalklasik}

\begin{solusi}
\textbf{(a)} 
\begin{itemize}[leftmargin=*]
    \item $(\implies)$ Misalkan $e_K$ bijektif (khususnya injektif). Andaikan $\gcd(a, 26) = d > 1$.
    Pilih $x_1 = 0$ dan $x_2 = \frac{26}{d}$. Karena $1 < d \le 26$, maka $0 < x_2 < 26$, sehingga $x_1 \not\equiv x_2 \pmod{26}$.
    Tinjau selisih nilai fungsi:
    \[
    e_K(x_2) - e_K(x_1) \equiv \left(a \cdot \frac{26}{d} + b\right) - (a \cdot 0 + b) = \frac{a}{d} \cdot 26 \equiv 0 \pmod{26},
    \]
    karena $d \mid a \implies \frac{a}{d} \in \mathbb{Z}$.
    Hal ini menghasilkan $e_K(x_2) \equiv e_K(x_1) \pmod{26}$ untuk $x_1 \not\equiv x_2$, kontradiksi dengan sifat injektif $e_K$.
    $\therefore \gcd(a, 26) = 1$.
    
    \item $(\impliedby)$ Misalkan $\gcd(a, 26) = 1$. Menurut Teorema Bézout, terdapat $a^{-1} \in \mathbb{Z}_{26}$ sedemikian hingga $a \cdot a^{-1} \equiv 1 \pmod{26}$.
    Tinjau persamaan $e_K(x_1) \equiv e_K(x_2) \pmod{26}$:
    \begin{align*}
    ax_1 + b &\equiv ax_2 + b \pmod{26} \\
    \implies ax_1 &\equiv ax_2 \pmod{26} \\
    \implies a^{-1}(ax_1) &\equiv a^{-1}(ax_2) \pmod{26} \\
    \implies 1 \cdot x_1 &\equiv 1 \cdot x_2 \pmod{26} \implies x_1 \equiv x_2 \pmod{26}.
    \end{align*}
    Maka $e_K$ injektif. Karena domain dan kodomain berhingga dan sama kardinalitasnya ($|\mathbb{Z}_{26}| = 26$), pemetaan injektif mengakibatkan surjektif, sehingga $e_K$ bijektif.
\end{itemize}

\textbf{(b)} Diberikan $y \equiv ax + b \pmod{26}$. Selesaikan untuk $x$:
\begin{align*}
y - b &\equiv ax \pmod{26} \\
\iff a^{-1}(y - b) &\equiv a^{-1}(ax) \pmod{26} \\
\iff x &\equiv a^{-1}(y - b) \pmod{26}.
\end{align*}
$\therefore d_K(y) \equiv a^{-1}(y - b) \pmod{26}$.

Buktikan komposisi fungsi:
\begin{align*}
d_K(e_K(x)) &\equiv a^{-1}\big((ax + b) - b\big) \pmod{26} \\
&\equiv a^{-1}(ax) \pmod{26} \\
&\equiv (a^{-1}a)x \pmod{26} \\
&\equiv 1 \cdot x \equiv x \pmod{26}.
\end{align*}
\end{solusi}

\begin{soalklasik}{Soal 3 [Klasik - Affine Modulo Prima]}
Diberikan sandi Affine atas ruang residu prima $\mathbb{Z}_{29}$ dengan kunci $K = (a, b) = (5, 21)$.
\begin{enumerate}[label=(\alph*)]
    \item Tentukan invers perkalian dari $a = 5$ dalam $\mathbb{Z}_{29}$, lalu nyatakan fungsi dekripsi dalam bentuk $d_K(y) \equiv cy + d \pmod{29}$ dengan $c, d \in \mathbb{Z}_{29}$.
    \item Dekripsi simbol ciphertext $y = 10 \in \mathbb{Z}_{29}$ menggunakan formula yang Anda peroleh.
\end{enumerate}
\end{soalklasik}

\begin{solusi}
\textbf{(a)} Menentukan $5^{-1} \pmod{29}$:
\[
5 \times 6 = 30 = 1 \times 29 + 1 \equiv 1 \pmod{29} \implies 5^{-1} \equiv 6 \pmod{29}.
\]
Fungsi dekripsi:
\begin{align*}
d_K(y) &\equiv 5^{-1}(y - 21) \pmod{29} \\
&\equiv 6(y - 21) \pmod{29} \\
&\equiv 6y - 126 \pmod{29}.
\end{align*}
Reduksi konstanta modulo 29:
\[
-126 \equiv -126 + 5 \times 29 = -126 + 145 = 19 \pmod{29}.
\]
$\therefore d_K(y) \equiv 6y + 19 \pmod{29}$, sehingga $c = 6$ dan $d = 19$.

\textbf{(b)} Untuk $y = 10$:
\begin{align*}
d_K(10) &\equiv 6(10) + 19 \pmod{29} \\
&\equiv 60 + 19 = 79 \pmod{29} \\
&\equiv 79 - 2 \times 29 = 79 - 58 = 21 \pmod{29}.
\end{align*}
$\therefore x = 21$.
\end{solusi}

\begin{soalklasik}{Soal 4 [Klasik - Hill Inversibilitas]}
Diberikan matriks kunci ukuran $2 \times 2$ untuk Sandi Hill:
\[
M = \begin{pmatrix} a & b \\ 2 & a \end{pmatrix} \in \mathcal{M}_{2 \times 2}(\mathbb{Z}_{26}).
\]
\begin{enumerate}[label=(\alph*)]
    \item Tentukan determinan $\det(M)$ sebagai fungsi dari $a$ dan $b$.
    \item Tentukan syarat perlu dan cukup pada pasangan $(a, b)$ agar matriks $M$ memiliki invers perkalian atas $\mathbb{Z}_{26}$.
    \item Buktikan bahwa nilai $a$ harus merupakan bilangan bulat ganjil dan $b \not\equiv 7a^2 \pmod{13}$.
\end{enumerate}
\end{soalklasik}

\begin{solusi}
\textbf{(a)} Determinan matriks $M$:
\[
\det(M) = a \cdot a - 2 \cdot b = a^2 - 2b \pmod{26}.
\]

\textbf{(b)} Matriks $M \in \mathcal{M}_{2 \times 2}(\mathbb{Z}_{26})$ memiliki invers $\iff \gcd(\det(M), 26) = 1$.

\textbf{(c)} Faktorisasi prima: $26 = 2 \times 13$. Syarat $\gcd(a^2 - 2b, 26) = 1$ ekuivalen dengan sistem:
\[
\begin{cases}
a^2 - 2b \not\equiv 0 \pmod 2 \\
a^2 - 2b \not\equiv 0 \pmod{13}
\end{cases}
\]
\begin{enumerate}[leftmargin=*]
    \item \textbf{Tinjauan Modulo 2:}
    \[
    a^2 - 2b \equiv a^2 - 0 \equiv a^2 \pmod 2.
    \]
    Maka:
    \[
    a^2 - 2b \not\equiv 0 \pmod 2 \iff a^2 \equiv 1 \pmod 2 \iff a \equiv 1 \pmod 2 \iff a \text{ ganjil}.
    \]
    
    \item \textbf{Tinjauan Modulo 13:}
    \[
    a^2 - 2b \not\equiv 0 \pmod{13} \iff 2b \not\equiv a^2 \pmod{13}.
    \]
    Karena $2 \times 7 = 14 \equiv 1 \pmod{13}$, maka $2^{-1} \equiv 7 \pmod{13}$. Kalikan kedua ruas dengan 7:
    \[
    b \not\equiv 7a^2 \pmod{13}.
    \]
\end{enumerate}
$\therefore M$ invertible atas $\mathbb{Z}_{26} \iff a \text{ ganjil}$ dan $b \not\equiv 7a^2 \pmod{13}$.
\end{solusi}

\begin{soalklasik}{Soal 5 [Klasik - Hill 2x2 Dekripsi]}
Diberikan kunci matriks Sandi Hill:
\[
K = \begin{pmatrix} 3 & 1 \\ 2 & 3 \end{pmatrix} \in \mathcal{M}_{2 \times 2}(\mathbb{Z}_{26}).
\]
\begin{enumerate}[label=(\alph*)]
    \item Tentukan determinan $K$, buktikan $K$ memiliki invers di $\mathbb{Z}_{26}$, dan hitung matriks dekripsi $K^{-1} \pmod{26}$.
    \item Lakukan dekripsi terhadap blok ciphertext $\mathbf{y} = \begin{pmatrix} 7 \\ 15 \end{pmatrix} \pmod{26}$.
\end{enumerate}
\end{soalklasik}

\begin{solusi}
\textbf{(a)} Determinan:
\[
\det(K) = 3(3) - 1(2) = 9 - 2 = 7.
\]
Karena $\gcd(7, 26) = 1$, maka $K$ invertible atas $\mathbb{Z}_{26}$.
Mencari $7^{-1} \pmod{26}$:
\[
7 \times 15 = 105 = 4 \times 26 + 1 \equiv 1 \pmod{26} \implies 7^{-1} \equiv 15 \pmod{26}.
\]
Adjoin matriks $K$:
\[
\operatorname{adj}(K) = \begin{pmatrix} 3 & -1 \\ -2 & 3 \end{pmatrix} \equiv \begin{pmatrix} 3 & 25 \\ 24 & 3 \end{pmatrix} \pmod{26}.
\]
Matriks invers $K^{-1}$:
\begin{align*}
K^{-1} &\equiv 15 \begin{pmatrix} 3 & -1 \\ -2 & 3 \end{pmatrix} = \begin{pmatrix} 45 & -15 \\ -30 & 45 \end{pmatrix} \pmod{26} \\
&\equiv \begin{pmatrix} 19 & 11 \\ 22 & 19 \end{pmatrix} \pmod{26}.
\end{align*}

\textbf{(b)} Dekripsi $\mathbf{y} = \begin{pmatrix} 7 \\ 15 \end{pmatrix}$:
\begin{align*}
\mathbf{x} &\equiv K^{-1} \mathbf{y} \pmod{26} \\
&= \begin{pmatrix} 19 & 11 \\ 22 & 19 \end{pmatrix} \begin{pmatrix} 7 \\ 15 \end{pmatrix} \\
&= \begin{pmatrix} 19(7) + 11(15) \\ 22(7) + 19(15) \end{pmatrix} = \begin{pmatrix} 133 + 165 \\ 154 + 285 \end{pmatrix} = \begin{pmatrix} 298 \\ 439 \end{pmatrix} \pmod{26}.
\end{align*}
Reduksi modulo 26:
\[
298 = 11 \times 26 + 12 \equiv 12 \pmod{26}, \quad 439 = 16 \times 26 + 23 \equiv 23 \pmod{26}.
\]
$\therefore \mathbf{x} = \begin{pmatrix} 12 \\ 23 \end{pmatrix} \pmod{26}$.
\end{solusi}

\begin{soalklasik}{Soal 6 [Klasik - Hill 3x3 Aljabar]}
Diberikan matriks $3 \times 3$ atas $\mathbb{Z}_{26}$:
\[
A = \begin{pmatrix} 1 & 0 & 1 \\ 0 & a & 1 \\ 1 & 2 & 1 \end{pmatrix}.
\]
\begin{enumerate}[label=(\alph*)]
    \item Hitung determinan matriks $A$ dalam $\mathbb{Z}_{26}$.
    \item Tentukan semua nilai parameter $a \in \mathbb{Z}_{26}$ yang memungkinkan matriks $A$ dapat digunakan sebagai matriks kunci enkripsi Sandi Hill. Jelaskan implikasi aljabar dari hasil yang Anda peroleh.
\end{enumerate}
\end{soalklasik}

\begin{solusi}
\textbf{(a)} Ekspansi kofaktor pada baris pertama:
\begin{align*}
\det(A) &= 1 \cdot \det\begin{pmatrix} a & 1 \\ 2 & 1 \end{pmatrix} - 0 \cdot \det\begin{pmatrix} 0 & 1 \\ 1 & 1 \end{pmatrix} + 1 \cdot \det\begin{pmatrix} 0 & a \\ 1 & 2 \end{pmatrix} \\
&= 1(a \cdot 1 - 1 \cdot 2) + 1(0 \cdot 2 - a \cdot 1) \\
&= (a - 2) - a = -2 \equiv 24 \pmod{26}.
\end{align*}

\textbf{(b)} Matriks $A$ invertible $\iff \gcd(\det(A), 26) = 1$.
Namun untuk setiap $a \in \mathbb{Z}_{26}$, $\det(A) \equiv 24 \pmod{26}$.
Evaluasi FPB:
\[
\gcd(24, 26) = 2 \ne 1.
\]
Karena $\gcd(\det(A), 26) = 2 \ne 1$ untuk setiap $a$, elemen $\det(A)$ adalah pembagi nol (\textit{zero divisor}) di $\mathbb{Z}_{26}$ dan tidak memiliki invers perkalian.
$\therefore$ Himpunan parameter $a$ yang memenuhi adalah $\emptyset$ (tidak ada).
\end{solusi}

\begin{soalklasik}{Soal 7 [Klasik - Transposisi Matriks n x m]}
Misalkan suatu pesan berpanjang $L = n \cdot m$ dituliskan ke dalam matriks berukuran $n \times m$ ($n$ baris, $m$ kolom) secara mendatar (baris demi baris), lalu ciphertext dibentuk dengan membaca matriks tersebut secara vertikal (kolom demi kolom dari kiri ke kanan).
\begin{enumerate}[label=(\alph*)]
    \item Tuliskan pemetaan analitis indeks $k = \pi(t)$ yang menyatakan posisi karakter ke-$t$ pada plaintext ($0 \le t < L$) menjadi posisi ke-$k$ pada ciphertext ($0 \le k < L$).
    \item Berikan formula analitis pemetaan balik (dekripsi) untuk merekonstruksi posisi plaintext dari indeks ciphertext.
    \item Dekripsi ciphertext berikut yang diketahui berukuran $n = 6$ dan $m = 5$:
    \[
    \texttt{"SIAAEAANLNRSNIADMIDAHIUYIDSPTA"}
    \]
\end{enumerate}
\end{soalklasik}

\begin{solusi}
\textbf{(a)} Karakter plaintext ke-$t \in \{0, 1, \dots, L-1\}$ menempati koordinat:
\[
\text{baris } i = \lfloor t / m \rfloor, \quad \text{kolom } j = t \bmod m.
\]
Pembacaan kolom demi kolom (setiap kolom berpanjang $n$):
\[
k = \pi(t) = j \cdot n + i = (t \bmod m) \cdot n + \lfloor t / m \rfloor.
\]

\textbf{(b)} Karakter ciphertext ke-$k \in \{0, 1, \dots, L-1\}$ berasal dari:
\[
\text{kolom } j = \lfloor k / n \rfloor, \quad \text{baris } i = k \bmod n.
\]
Maka indeks posisi karakter pada plaintext adalah:
\[
t = \pi^{-1}(k) = i \cdot m + j = (k \bmod n) \cdot m + \lfloor k / n \rfloor.
\]

\textbf{(c)} Diberikan $L = 30$, $n = 6$ (baris), $m = 5$ (kolom).
Ciphertext diisikan per kolom berukuran $n = 6$:
\[
\begin{array}{c|ccccc}
\text{Baris} & \text{Kolom } 0 & \text{Kolom } 1 & \text{Kolom } 2 & \text{Kolom } 3 & \text{Kolom } 4 \\ \hline
0 & \texttt{S} & \texttt{A} & \texttt{N} & \texttt{D} & \texttt{I} \\
1 & \texttt{I} & \texttt{N} & \texttt{I} & \texttt{A} & \texttt{D} \\
2 & \texttt{A} & \texttt{L} & \texttt{A} & \texttt{H} & \texttt{S} \\
3 & \texttt{A} & \texttt{N} & \texttt{D} & \texttt{I} & \texttt{P} \\
4 & \texttt{E} & \texttt{R} & \texttt{M} & \texttt{U} & \texttt{T} \\
5 & \texttt{A} & \texttt{S} & \texttt{I} & \texttt{Y} & \texttt{A}
\end{array}
\]
Membaca baris demi baris:
\[
\texttt{"SANDI INI ADALAH SANDI PERMUTASI YA"}.
\]
\end{solusi}

\begin{soalklasik}{Soal 8 [Klasik - Vigenère Komposisi]}
Misalkan suatu teks-asal dienkripsi berturut-turut menggunakan dua buah sandi Vigenère dengan panjang kunci masing-masing $m_1$ dan $m_2$.
Buktikan secara formal bahwa sistem enkripsi komposit tersebut setara dengan sebuah sandi Vigenère tunggal dengan panjang kunci paling banyak $\operatorname{kpk}(m_1, m_2)$.
\end{soalklasik}

\begin{solusi}
Misalkan $\mathbf{k}_1 = (k_{1, 0}, \dots, k_{1, m_1-1}) \in \mathbb{Z}_{26}^{m_1}$ dan $\mathbf{k}_2 = (k_{2, 0}, \dots, k_{2, m_2-1}) \in \mathbb{Z}_{26}^{m_2}$.
Enkripsi tahap 1:
\[
y_i \equiv p_i + k_{1, i \bmod m_1} \pmod{26}.
\]
Enkripsi tahap 2:
\[
c_i \equiv y_i + k_{2, i \bmod m_2} \equiv p_i + \big(k_{1, i \bmod m_1} + k_{2, i \bmod m_2}\big) \pmod{26}.
\]
Definisikan barisan pergeseran komposit $s_i \in \mathbb{Z}_{26}$:
\[
s_i \equiv k_{1, i \bmod m_1} + k_{2, i \bmod m_2} \pmod{26}.
\]
Misalkan $M = \operatorname{kpk}(m_1, m_2)$. Maka $m_1 \mid M$ dan $m_2 \mid M$.
Untuk setiap $i \ge 0$:
\begin{align*}
s_{i + M} &\equiv k_{1, (i + M) \bmod m_1} + k_{2, (i + M) \bmod m_2} \pmod{26} \\
&\equiv k_{1, i \bmod m_1} + k_{2, i \bmod m_2} \pmod{26} \\
&\equiv s_i \pmod{26}.
\end{align*}
Barisan $(s_i)$ periodik dengan periode yang membagi $M$.
Definisikan vektor kunci tunggal $\mathbf{k}_3 = (s_0, s_1, \dots, s_{M-1}) \in \mathbb{Z}_{26}^M$.
Maka $c_i \equiv p_i + k_{3, i \bmod M} \pmod{26}$.
$\therefore$ Sistem komposit ekuivalen dengan sandi Vigenère berpanjang kunci paling banyak $\operatorname{kpk}(m_1, m_2)$.
\end{solusi}

\newpage
% =========================================================================
\section{Solusi Bagian II: Kriptoanalisis (Soal 9--13)}
% =========================================================================

\begin{soalanalisis}{Soal 9 [Kriptoanalisis - Kasiski]}
Diberikan ciphertext hasil enkripsi Vigenère sebagai berikut:
\[
\texttt{KXQMYTOZKQWHWMEKXQQXKQXYM}
\]
\begin{enumerate}[label=(\alph*)]
    \item Identifikasi seluruh posisi kemunculan substring berulang $\texttt{KXQ}$ dan $\texttt{KQ}$ (indeks berbasis 1).
    \item Hitung seluruh jarak selisih antar posisi perulangan tersebut, faktorkan ke dalam perkalian bilangan prima, dan lakukan pengujian Kasiski untuk menduga panjang kunci $m$.
\end{enumerate}
\end{soalanalisis}

\begin{solusi}
\textbf{(a)} Ciphertext berpanjang $N = 25$:
\[
\begin{array}{l*{13}{c}}
\text{Indeks:} & 1 & 2 & 3 & 4 & 5 & 6 & 7 & 8 & 9 & 10 & 11 & 12 & 13 \\
\text{Karakter:} & \texttt{K} & \texttt{X} & \texttt{Q} & \texttt{M} & \texttt{Y} & \texttt{T} & \texttt{O} & \texttt{Z} & \texttt{K} & \texttt{Q} & \texttt{W} & \texttt{H} & \texttt{W} \\[0.5em]
\text{Indeks:} & 14 & 15 & 16 & 17 & 18 & 19 & 20 & 21 & 22 & 23 & 24 & 25 & \\
\text{Karakter:} & \texttt{M} & \texttt{E} & \texttt{K} & \texttt{X} & \texttt{Q} & \texttt{Q} & \texttt{X} & \texttt{K} & \texttt{Q} & \texttt{X} & \texttt{Y} & \texttt{M} & \\
\end{array}
\]
Identifikasi posisi:
\begin{itemize}
    \item Substring $\texttt{KXQ}$: posisi $1$ dan $16$.
    \item Substring $\texttt{KQ}$: posisi $9$ dan $21$.
\end{itemize}

\textbf{(b)} Perhitungan jarak dan faktorisasi:
\begin{itemize}
    \item Jarak $\texttt{KXQ}$: $\Delta_1 = 16 - 1 = 15 = 3 \times 5$.
    \item Jarak $\texttt{KQ}$: $\Delta_2 = 21 - 9 = 12 = 2^2 \times 3$.
    \item Jarak silang awalan $\texttt{K}$ pada pola:
    \begin{align*}
    21 - 1 &= 20 = 2^2 \times 5, \\
    16 - 9 &= 7 \text{ (prima)}.
    \end{align*}
\end{itemize}
Tinjau FPB jarak perulangan substring trigram $\texttt{KXQ}$ dan jarak kemunculan $\texttt{K}$ dominan:
\[
\gcd(15, 20) = 5.
\]
Karena trigram $\texttt{KXQ}$ berulang pada jarak $\Delta_1 = 15 = 3 \times 5$, maka kandidat panjang kunci menurut uji Kasiski adalah faktor dari 15, yaitu $m \in \{3, 5\}$. Munculnya kelipatan 5 pada jarak terjauh ($21 - 1 = 20$) memberikan dugaan terkuat:
\[
m = 5.
\]
\end{solusi}

\begin{soalanalisis}{Soal 10 [Kriptoanalisis - Indeks Koinsidensi]}
Diberikan teks sandi $\mathbf{x} = x_1 x_2 \dots x_N$ dengan panjang $N$ dan frekuensi kemunculan huruf ke-$i$ adalah $f_i$ ($0 \le i \le 25$).
\begin{enumerate}[label=(\alph*)]
    \item Buktikan bahwa nilai harapan Indeks Koinsidensi untuk string yang dibangkitkan dari distribusi peluang seragam $p_i = \frac{1}{26}$ adalah $E[I_c] = \frac{1}{26} \approx 0.03846$.
    \item Suatu ciphertext berukuran $N = 100$ memiliki tabulasi frekuensi tertentu yang menghasilkan $\sum_{i=0}^{25} f_i(f_i - 1) = 654$. Hitung nilai $I_c$ dan tentukan apakah ciphertext tersebut dienkripsi menggunakan cipher monoalfabetik atau polialfabetik.
\end{enumerate}
\end{soalanalisis}

\begin{solusi}
\textbf{(a)} Indeks Koinsidensi:
\[
I_c = \frac{\sum_{i=0}^{25} f_i(f_i - 1)}{N(N - 1)}.
\]
Tinjau variabel acak indikator untuk pasangan indeks $1 \le j < k \le N$:
\[
I_{j, k} = \begin{cases} 1, & \text{jika } x_j = x_k \\ 0, & \text{jika } x_j \ne x_k \end{cases}
\]
Maka $\sum_{i=0}^{25} \binom{f_i}{2} = \sum_{1 \le j < k \le N} I_{j, k}$. Karena setiap karakter dipilih independen dan seragam dengan $p_i = \frac{1}{26}$:
\[
P(x_j = x_k) = \sum_{i=0}^{25} P(x_j = i, x_k = i) = \sum_{i=0}^{25} p_i^2 = \sum_{i=0}^{25} \left(\frac{1}{26}\right)^2 = 26 \times \frac{1}{26^2} = \frac{1}{26}.
\]
Berdasarkan kelinieran nilai harapan:
\[
E\left[\sum_{i=0}^{25} f_i(f_i - 1)\right] = 2 E\left[\sum_{j < k} I_{j, k}\right] = 2 \binom{N}{2} \cdot \frac{1}{26} = N(N - 1) \cdot \frac{1}{26}.
\]
Maka:
\[
E[I_c] = \frac{N(N - 1) \cdot \frac{1}{26}}{N(N - 1)} = \frac{1}{26} \approx 0.03846.
\]

\textbf{(b)} Untuk $N = 100$ dan $\sum_{i=0}^{25} f_i(f_i - 1) = 654$:
\[
I_c = \frac{654}{100(99)} = \frac{654}{9900} \approx 0.06606.
\]
Analisis:
\begin{itemize}
    \item $I_{\text{mono}} \approx 0.065$ (teks alami monoalfabetik).
    \item $I_{\text{acak}} \approx 0.0385$ (teks acak / polialfabetik dengan kunci panjang).
\end{itemize}
Karena nilai $I_c = 0.0661$ sangat dekat dengan $0.065$, distribusi frekuensi huruf belum terseragamkan.
$\therefore$ Ciphertext dienkripsi menggunakan \textbf{sandi monoalfabetik}.
\end{solusi}

\begin{soalanalisis}{Soal 11 [Kriptoanalisis - Estimasi Panjang Kunci Vigenère]}
Suatu teks sandi Vigenère berpanjang $N = 500$ karakter menghasilkan Indeks Koinsidensi terhitung sebesar $I_c = 0.0450$.
Gunakan aproksimasi formula Babbage-Friedman:
\[
m \approx \frac{0.027 N}{(N - 1)I_c - 0.038 N + 0.065}
\]
untuk menaksir panjang kunci Vigenère $m$ yang digunakan.
\end{soalanalisis}

\begin{solusi}
Substitusikan parameter $N = 500$ dan $I_c = 0.0450$:
\begin{align*}
\text{Pembilang} &= 0.027 \times 500 = 13.5, \\
\text{Penyebut} &= (500 - 1)(0.0450) - 0.038(500) + 0.065 \\
&= 499 \times 0.0450 - 19.0 + 0.065 \\
&= 22.455 - 19.0 + 0.065 \\
&= 3.52.
\end{align*}
Maka:
\[
m \approx \frac{13.5}{3.52} \approx 3.835.
\]
Karena panjang kunci $m \in \mathbb{Z}^+$, pembulatan terdekat memberikan taksiran:
\[
m \approx 4.
\]
\end{solusi}

\begin{soalanalisis}{Soal 12 [Kriptoanalisis - Affine KPA]}
Seorang kriptanalis berhasil menyadap pasangan plaintext-ciphertext dari sandi Affine atas $\mathbb{Z}_{26}$:
\[
\texttt{'IF'} \longrightarrow \texttt{'PQ'}
\]
dengan korespondensi alfabet baku $\texttt{A}=0, \texttt{B}=1, \dots, \texttt{Z}=25$.
\begin{enumerate}[label=(\alph*)]
    \item Susun sistem kongruensi linier dalam $\mathbb{Z}_{26}$ untuk menentukan kunci $a$ dan $b$.
    \item Selesaikan sistem tersebut untuk memperoleh nilai eksak pasangan kunci $K = (a, b)$.
\end{enumerate}
\end{soalanalisis}

\begin{solusi}
\textbf{(a)} Konversi nilai numerik karakter:
\[
\texttt{I} = 8, \quad \texttt{F} = 5, \quad \texttt{P} = 15, \quad \texttt{Q} = 16.
\]
Sistem kongruensi dari $e_K(x) \equiv ax + b \pmod{26}$:
\begin{align}
8a + b &\equiv 15 \pmod{26} \\
5a + b &\equiv 16 \pmod{26}
\end{align}

\textbf{(b)} Kurangkan persamaan (1) dengan persamaan (2):
\[
(8a + b) - (5a + b) \equiv 15 - 16 \pmod{26} \implies 3a \equiv -1 \equiv 25 \pmod{26}.
\]
Invers perkalian dari $3$ modulo $26$:
\[
3 \times 9 = 27 \equiv 1 \pmod{26} \implies 3^{-1} \equiv 9 \pmod{26}.
\]
Kalikan kedua ruas dengan 9:
\[
a \equiv 9 \times 25 \equiv 9(-1) = -9 \equiv 17 \pmod{26}.
\]
Verifikasi: $\gcd(17, 26) = 1$.
Substitusikan $a = 17$ ke persamaan (2):
\begin{align*}
5(17) + b &\equiv 16 \pmod{26} \\
85 + b &\equiv 16 \pmod{26} \\
7 + b &\equiv 16 \pmod{26} \implies b \equiv 9 \pmod{26}.
\end{align*}
$\therefore$ Pasangan kunci adalah $K = (a, b) = (17, 9)$.
\end{solusi}

\begin{soalanalisis}{Soal 13 [Kriptoanalisis - Hill KPA 2x2]}
Melalui serangan \textit{Known-Plaintext Attack} pada Sandi Hill $2 \times 2$ atas $\mathbb{Z}_{26}$, diperoleh dua pasangan blok plaintext dan ciphertext:
\[
\mathbf{p}_1 = \begin{pmatrix} 3 \\ 5 \end{pmatrix} \to \mathbf{c}_1 = \begin{pmatrix} 19 \\ 24 \end{pmatrix}, \quad \mathbf{p}_2 = \begin{pmatrix} 2 \\ 7 \end{pmatrix} \to \mathbf{c}_2 = \begin{pmatrix} 20 \\ 7 \end{pmatrix}.
\]
\begin{enumerate}[label=(\alph*)]
    \item Bentuk persamaan matriks $C \equiv K P \pmod{26}$, di mana $P = [\mathbf{p}_1 \;\; \mathbf{p}_2]$ dan $C = [\mathbf{c}_1 \;\; \mathbf{c}_2]$.
    \item Buktikan bahwa matriks plaintext $P$ memiliki invers dalam $\mathbb{Z}_{26}$ dan tentukan matriks $P^{-1} \pmod{26}$.
    \item Hitung matriks kunci rahasia $K = C P^{-1} \pmod{26}$.
\end{enumerate}
\end{soalanalisis}

\begin{solusi}
\textbf{(a)} Susun matriks:
\[
P = \begin{pmatrix} 3 & 2 \\ 5 & 7 \end{pmatrix}, \quad C = \begin{pmatrix} 19 & 20 \\ 24 & 7 \end{pmatrix} \implies C \equiv K P \pmod{26}.
\]

\textbf{(b)} Determinan matriks $P$:
\[
\det(P) = 3(7) - 2(5) = 21 - 10 = 11.
\]
Karena $\gcd(11, 26) = 1$, maka $P$ invertible di $\mathbb{Z}_{26}$.
Invers $11$ modulo $26$:
\[
11 \times 19 = 209 = 8 \times 26 + 1 \equiv 1 \pmod{26} \implies 11^{-1} \equiv 19 \pmod{26}.
\]
Adjoin $P$:
\[
\operatorname{adj}(P) = \begin{pmatrix} 7 & -2 \\ -5 & 3 \end{pmatrix} \equiv \begin{pmatrix} 7 & 24 \\ 21 & 3 \end{pmatrix} \pmod{26}.
\]
Invers $P^{-1}$:
\begin{align*}
P^{-1} &\equiv 19 \begin{pmatrix} 7 & -2 \\ -5 & 3 \end{pmatrix} = \begin{pmatrix} 133 & -38 \\ -95 & 57 \end{pmatrix} \pmod{26} \\
&\equiv \begin{pmatrix} 3 & 14 \\ 9 & 5 \end{pmatrix} \pmod{26}.
\end{align*}

\textbf{(c)} Menghitung $K = C P^{-1} \pmod{26}$:
\begin{align*}
K &\equiv \begin{pmatrix} 19 & 20 \\ 24 & 7 \end{pmatrix} \begin{pmatrix} 3 & 14 \\ 9 & 5 \end{pmatrix} \pmod{26} \\
&= \begin{pmatrix} 19(3) + 20(9) & 19(14) + 20(5) \\ 24(3) + 7(9) & 24(14) + 7(5) \end{pmatrix} \\
&= \begin{pmatrix} 57 + 180 & 266 + 100 \\ 72 + 63 & 336 + 35 \end{pmatrix} = \begin{pmatrix} 237 & 366 \\ 135 & 371 \end{pmatrix} \pmod{26}.
\end{align*}
Reduksi modulo 26:
\begin{align*}
237 &= 9 \times 26 + 3 \equiv 3, & 366 &= 14 \times 26 + 2 \equiv 2, \\
135 &= 5 \times 26 + 5 \equiv 5, & 371 &= 14 \times 26 + 7 \equiv 7.
\end{align*}
$\therefore K = \begin{pmatrix} 3 & 2 \\ 5 & 7 \end{pmatrix}$.
\end{solusi}

\newpage
% =========================================================================
\section{Solusi Bagian III: Teori Shannon (Soal 14--20)}
% =========================================================================

\begin{soalshannon}{Soal 14 [Shannon - Perfect Secrecy Geser]}
Pandang sandi geser dengan $\mathcal{P} = \mathcal{C} = \mathcal{K} = \mathbb{Z}_{26}$. Kunci $K$ dipilih secara acak seragam dengan peluang $p(K = k) = \frac{1}{26}$ untuk setiap $k \in \mathbb{Z}_{26}$, dan pemilihan kunci saling bebas terhadap plaintext $P$.
Buktikan secara formal bahwa sistem sandi geser tersebut memenuhi definisi keamanan sempurna (\textit{perfect secrecy}):
\[
p(x \mid y) = p(x), \quad \forall x \in \mathbb{Z}_{26}, \forall y \in \mathbb{Z}_{26}
\]
untuk sebarang distribusi probabilitas a priori $p(x)$ pada $\mathcal{P}$.
\end{soalshannon}

\begin{solusi}
Untuk setiap $y \in \mathbb{Z}_{26}$, probabilitas marginal ciphertext adalah:
\[
p(y) = \sum_{x' \in \mathbb{Z}_{26}} P(P = x', C = y) = \sum_{x' \in \mathbb{Z}_{26}} p(x') P(K = (y - x') \bmod 26).
\]
Karena $K$ berdistribusi seragam dan independen terhadap $P$, $P(K = (y - x') \bmod 26) = \frac{1}{26}$ untuk setiap $x'$. Maka:
\[
p(y) = \sum_{x' \in \mathbb{Z}_{26}} p(x') \cdot \frac{1}{26} = \frac{1}{26} \sum_{x' \in \mathbb{Z}_{26}} p(x') = \frac{1}{26} \cdot 1 = \frac{1}{26}.
\]
Berdasarkan Definisi Peluang Bersyarat (Aturan Bayes):
\begin{align*}
p(x \mid y) &= \frac{P(P = x, C = y)}{p(y)} \\
&= \frac{p(x) P(K = (y - x) \bmod 26)}{p(y)} \\
&= \frac{p(x) \cdot \frac{1}{26}}{\frac{1}{26}} = p(x).
\end{align*}
Karena $p(x \mid y) = p(x)$ berlaku untuk setiap $x \in \mathcal{P}$ dan setiap $y \in \mathcal{C}$, maka sistem memenuhi definisi keamanan sempurna (\textit{perfect secrecy}).
\end{solusi}

\begin{soalshannon}{Soal 15 [Shannon - Matriks Probabilitas C]}
Pandang sistem kripto dengan ruang pesan $\mathcal{P} = \{a, b, c\}$, ruang kunci $\mathcal{K} = \{K_1, K_2, K_3\}$, dan ruang teks sandi $\mathcal{C} = \{1, 2, 3, 4\}$. Matriks enkripsi didefinisikan sebagai:
\[
\begin{array}{c|ccc}
e_K(p) & a & b & c \\ \hline
K_1 & 1 & 2 & 3 \\
K_2 & 2 & 3 & 4 \\
K_3 & 3 & 4 & 1
\end{array}
\]
Distribusi peluang pemilihan kunci adalah seragam: $p(K_1) = p(K_2) = p(K_3) = \frac{1}{3}$.
Distribusi peluang teks-asal diberikan oleh: $p(a) = \frac{1}{2}, p(b) = \frac{1}{3}, p(c) = \frac{1}{6}$.
Hitung nilai eksak distribusi probabilitas marginal ciphertext:
\[
p(C = 1), \quad p(C = 2), \quad p(C = 3), \quad \text{dan} \quad p(C = 4).
\]
\end{soalshannon}

\begin{solusi}
Formula marginal: $p(y) = \sum_{x \in \mathcal{P}} \sum_{\{k : e_k(x) = y\}} p(x) p(k) = \frac{1}{3} \sum_{\{x : \exists k, e_k(x)=y\}} p(x)$.
\begin{itemize}[leftmargin=*]
    \item Untuk $y = 1$: dihasilkan oleh pasangan $(K_1, a)$ dan $(K_3, c)$.
    \[
    p(1) = \frac{1}{3} p(a) + \frac{1}{3} p(c) = \frac{1}{3}\left(\frac{1}{2} + \frac{1}{6}\right) = \frac{1}{3}\left(\frac{4}{6}\right) = \frac{4}{18} = \frac{2}{9}.
    \]
    
    \item Untuk $y = 2$: dihasilkan oleh pasangan $(K_1, b)$ dan $(K_2, a)$.
    \[
    p(2) = \frac{1}{3} p(b) + \frac{1}{3} p(a) = \frac{1}{3}\left(\frac{1}{3} + \frac{1}{2}\right) = \frac{1}{3}\left(\frac{5}{6}\right) = \frac{5}{18}.
    \]
    
    \item Untuk $y = 3$: dihasilkan oleh pasangan $(K_1, c)$, $(K_2, b)$, dan $(K_3, a)$.
    \[
    p(3) = \frac{1}{3} p(c) + \frac{1}{3} p(b) + \frac{1}{3} p(a) = \frac{1}{3}\left(\frac{1}{6} + \frac{1}{3} + \frac{1}{2}\right) = \frac{1}{3}(1) = \frac{1}{3} = \frac{6}{18}.
    \]
    
    \item Untuk $y = 4$: dihasilkan oleh pasangan $(K_2, c)$ dan $(K_3, b)$.
    \[
    p(4) = \frac{1}{3} p(c) + \frac{1}{3} p(b) = \frac{1}{3}\left(\frac{1}{6} + \frac{1}{3}\right) = \frac{1}{3}\left(\frac{3}{6}\right) = \frac{3}{18} = \frac{1}{6}.
    \]
\end{itemize}
Verifikasi total probabilitas:
\[
\sum_{y=1}^4 p(y) = \frac{4}{18} + \frac{5}{18} + \frac{6}{18} + \frac{3}{18} = \frac{18}{18} = 1.
\]
$\therefore p(1) = \frac{2}{9}, \; p(2) = \frac{5}{18}, \; p(3) = \frac{1}{3}, \; p(4) = \frac{1}{6}$.
\end{solusi}

\begin{soalshannon}{Soal 16 [Shannon - Entropi H(P), H(K), H(C)]}
Berdasarkan sistem kripto pada Soal 15:
\begin{enumerate}[label=(\alph*)]
    \item Hitung nilai entropi teks-asal $H(P)$ dan entropi kunci $H(K)$ dalam satuan bit (berikan representasi analitis eksak dan hampiran desimal 4 angka di belakang koma).
    \item Hitung nilai entropi ciphertext $H(C)$ dalam satuan bit.
\end{enumerate}
\end{soalshannon}

\begin{solusi}
\textbf{(a)} Entropi plaintext $H(P)$ dengan $p(a) = 1/2, p(b) = 1/3, p(c) = 1/6$:
\begin{align*}
H(P) &= -\left(\frac{1}{2}\log_2\frac{1}{2} + \frac{1}{3}\log_2\frac{1}{3} + \frac{1}{6}\log_2\frac{1}{6}\right) \\
&= \frac{1}{2}(1) + \frac{1}{3}\log_2 3 + \frac{1}{6}(\log_2 2 + \log_2 3) \\
&= \frac{1}{2} + \frac{1}{6} + \left(\frac{1}{3} + \frac{1}{6}\right)\log_2 3 \\
&= \frac{2}{3} + \frac{1}{2}\log_2 3.
\end{align*}
Nilai numerik ($\log_2 3 \approx 1.58496$):
\[
H(P) \approx 0.6667 + 0.5(1.58496) = 1.4591 \text{ bit}.
\]
Entropi kunci seragam $|\mathcal{K}| = 3$:
\[
H(K) = \log_2 3 \approx 1.5850 \text{ bit}.
\]

\textbf{(b)} Entropi ciphertext $H(C)$ dengan $p = \left(\frac{4}{18}, \frac{5}{18}, \frac{6}{18}, \frac{3}{18}\right)$:
\begin{align*}
H(C) &= -\sum_{y=1}^4 p(y) \log_2 p(y) \\
&= -\left[\frac{4}{18}\log_2\left(\frac{4}{18}\right) + \frac{5}{18}\log_2\left(\frac{5}{18}\right) + \frac{6}{18}\log_2\left(\frac{6}{18}\right) + \frac{3}{18}\log_2\left(\frac{3}{18}\right)\right] \\
&= \log_2 18 - \frac{1}{18}\big(4\log_2 4 + 5\log_2 5 + 6\log_2 6 + 3\log_2 3\big) \\
&= (1 + 2\log_2 3) - \frac{1}{18}\big(4(2) + 5\log_2 5 + 6(1 + \log_2 3) + 3\log_2 3\big) \\
&= 1 + 2\log_2 3 - \frac{1}{18}\big(14 + 5\log_2 5 + 9\log_2 3\big).
\end{align*}
Evaluasi numerik:
\begin{align*}
-\frac{4}{18}\log_2\left(\frac{4}{18}\right) &= -\frac{2}{9}\log_2\left(\frac{2}{9}\right) \approx 0.4822, \\
-\frac{5}{18}\log_2\left(\frac{5}{18}\right) &\approx 0.5133, \\
-\frac{6}{18}\log_2\left(\frac{6}{18}\right) &= -\frac{1}{3}\log_2\left(\frac{1}{3}\right) \approx 0.5283, \\
-\frac{3}{18}\log_2\left(\frac{3}{18}\right) &= -\frac{1}{6}\log_2\left(\frac{1}{6}\right) \approx 0.4308.
\end{align*}
Jumlahkan:
\[
H(C) \approx 0.4822 + 0.5133 + 0.5283 + 0.4308 = 1.9547 \text{ bit}.
\]
\end{solusi}

\begin{soalshannon}{Soal 17 [Shannon - Equivocation H(K|C) dan H(P|C)]}
Berdasarkan sistem kripto pada Soal 15:
\begin{enumerate}[label=(\alph*)]
    \item Buktikan secara deduktif penurunan identitas \textit{Key Equivocation}:
    \[
    H(K \mid C) = H(K) + H(P) - H(C)
    \]
    untuk sistem kripto dengan dekripsi deterministik dan kunci yang saling bebas terhadap pesan.
    \item Hitung nilai numerik $H(K \mid C)$ dan $H(P \mid C)$.
    \item Berdasarkan nilai $H(P \mid C)$ dan distribusi bersyarat $p(a \mid 1)$, buktikan apakah sistem ini memenuhi atau melanggar syarat keamanan sempurna (\textit{perfect secrecy}).
\end{enumerate}
\end{soalshannon}

\begin{solusi}
\textbf{(a)} Ekspansikan entropi gabungan $H(P, K, C)$ melalui dua urutan pengkondisian aturan rantai:
\begin{enumerate}[leftmargin=*]
    \item Urutan 1: $H(P, K, C) = H(P, K) + H(C \mid P, K)$.
    \begin{itemize}
        \item Karena $P$ dan $K$ saling bebas: $H(P, K) = H(P) + H(K)$.
        \item Karena enkripsi bersifat deterministik ($C = e_K(P)$), nilai $C$ pasti jika $P$ dan $K$ diketahui, sehingga $H(C \mid P, K) = 0$.
    \end{itemize}
    Maka:
    \[
    H(P, K, C) = H(P) + H(K).
    \]
    
    \item Urutan 2: $H(P, K, C) = H(C) + H(K \mid C) + H(P \mid K, C)$.
    \begin{itemize}
        \item Karena dekripsi bersifat deterministik ($P = d_K(C)$), nilai $P$ pasti jika $K$ dan $C$ diketahui, sehingga $H(P \mid K, C) = 0$.
    \end{itemize}
    Maka:
    \[
    H(P, K, C) = H(C) + H(K \mid C).
    \]
\end{enumerate}
Samakan kedua persamaan:
\[
H(C) + H(K \mid C) = H(P) + H(K) \implies H(K \mid C) = H(K) + H(P) - H(C).
\]

\textbf{(b)} Perhitungan numerik $H(K \mid C)$:
\[
H(K \mid C) = 1.5850 + 1.4591 - 1.9547 = 1.0894 \text{ bit}.
\]
Perhitungan $H(P \mid C)$:
Perhatikan bahwa pada matriks enkripsi Soal 15, setiap kolom memuat simbol ciphertext yang berbeda, dan setiap baris memuat simbol ciphertext yang berbeda. Artinya, untuk suatu $C = y$ yang teramati, mengetahui pesan $p$ menentukan secara unik kunci $k$ yang digunakan ($H(K \mid P, C) = 0$).
Akibatnya:
\[
H(P, K \mid C) = H(K \mid C) + H(P \mid K, C) = H(K \mid C) + 0 = H(K \mid C),
\]
dan di sisi lain:
\[
H(P, K \mid C) = H(P \mid C) + H(K \mid P, C) = H(P \mid C) + 0 = H(P \mid C).
\]
$\therefore H(P \mid C) = H(K \mid C) = 1.0894 \text{ bit}$.

\textbf{(c)} Tinjau syarat keamanan sempurna:
\begin{itemize}
    \item Ditinjau dari entropi: Sistem aman sempurna $\iff H(P \mid C) = H(P)$.
    Diperoleh:
    \[
    H(P \mid C) = 1.0894 \ne 1.4591 = H(P).
    \]
    \item Ditinjau dari probabilitas bersyarat $p(a \mid 1)$:
    \[
    p(a \mid 1) = \frac{P(P=a, C=1)}{p(C=1)} = \frac{p(a) p(K_1)}{p(1)} = \frac{\frac{1}{2} \cdot \frac{1}{3}}{\frac{4}{18}} = \frac{\frac{1}{6}}{\frac{4}{18}} = \frac{3}{4} = 0.75.
    \]
    Sedangkan $p(a) = 0.50$. Karena $p(a \mid 1) \ne p(a)$, mengamati ciphertext $C=1$ mengubah keyakinan tentang pesan $a$.
\end{itemize}
$\therefore$ Sistem \textbf{tidak memiliki keamanan sempurna}.
\end{solusi}

\begin{soalklasik}{Soal 18 [Shannon - Teorema Karakterisasi]}
Buktikan Teorema Shannon berikut:
Andaikan suatu sistem kripto $(\mathcal{P}, \mathcal{C}, \mathcal{K}, \mathcal{E}, \mathcal{D})$ memenuhi $|\mathcal{K}| = |\mathcal{C}| = |\mathcal{P}|$.
Sistem tersebut memiliki keamanan sempurna jika dan hanya jika:
\begin{enumerate}[label=(\arabic*)]
    \item Setiap kunci $k \in \mathcal{K}$ digunakan dengan probabilitas seragam: $p(k) = \frac{1}{|\mathcal{K}|}$.
    \item Untuk setiap $x \in \mathcal{P}$ dan setiap $y \in \mathcal{C}$, terdapat tepat satu kunci $k \in \mathcal{K}$ sedemikian hingga $e_k(x) = y$.
\end{enumerate}
\end{soalklasik}

\begin{solusi}
\textbf{Bukti $(\impliedby)$:}
Misalkan syarat (1) dan (2) dipenuhi.
Untuk setiap $x \in \mathcal{P}$ dan $y \in \mathcal{C}$, terdapat tepat satu kunci $k_{x, y} \in \mathcal{K}$ dengan $e_{k_{x, y}}(x) = y$.
Distribusi marginal $p(y)$:
\[
p(y) = \sum_{x' \in \mathcal{P}} p(x') P(e_K(x') = y) = \sum_{x' \in \mathcal{P}} p(x') p(k_{x', y}) = \frac{1}{|\mathcal{K}|} \sum_{x' \in \mathcal{P}} p(x') = \frac{1}{|\mathcal{K}|}.
\]
Probabilitas bersyarat $p(x \mid y)$:
\[
p(x \mid y) = \frac{P(P=x, C=y)}{p(y)} = \frac{p(x) p(k_{x, y})}{p(y)} = \frac{p(x) \cdot \frac{1}{|\mathcal{K}|}}{\frac{1}{|\mathcal{K}|}} = p(x).
\]
$\therefore$ Sistem aman sempurna.

\textbf{Bukti $(\implies)$:}
Misalkan sistem memiliki keamanan sempurna: $p(x \mid y) = p(x)$ untuk setiap $x \in \mathcal{P}, y \in \mathcal{C}$.
Asumsikan tanpa mengurangi keumuman bahwa $p(x) > 0$ untuk setiap $x \in \mathcal{P}$.
Dari definisi keamanan sempurna:
\[
p(y \mid x) = \frac{p(x \mid y) p(y)}{p(x)} = \frac{p(x) p(y)}{p(x)} = p(y) > 0.
\]
Karena $p(y \mid x) > 0$, maka untuk setiap $x \in \mathcal{P}$ dan setiap $y \in \mathcal{C}$, harus terdapat sedikitnya satu kunci $k \in \mathcal{K}$ sedemikian hingga $e_k(x) = y$.
Definisikan himpunan kunci $K(x, y) = \{k \in \mathcal{K} : e_k(x) = y\}$. Maka:
\[
|K(x, y)| \ge 1, \quad \forall x \in \mathcal{P}, \forall y \in \mathcal{C}.
\]
Untuk suatu $x \in \mathcal{P}$ yang tetap, karena enkripsi adalah pemetaan fungsi ($e_k(x)$ bernilai tunggal), himpunan-himpunan $\{K(x, y) : y \in \mathcal{C}\}$ saling asing dan membentuk partisi dari $\mathcal{K}$:
\[
\bigcup_{y \in \mathcal{C}} K(x, y) = \mathcal{K}.
\]
Akibatnya:
\[
|\mathcal{K}| = \sum_{y \in \mathcal{C}} |K(x, y)| \ge \sum_{y \in \mathcal{C}} 1 = |\mathcal{C}|.
\]
Karena diketahui $|\mathcal{K}| = |\mathcal{C}|$, ketaksamaan di atas haruslah menjadi kesamaan murni. Hal ini mengharuskan:
\[
|K(x, y)| = 1, \quad \forall x \in \mathcal{P}, \forall y \in \mathcal{C}.
\]
Ini membuktikan syarat (2): terdapat tepat satu kunci $k$ sedemikian hingga $e_k(x) = y$.

Selanjutnya, karena $|K(x, y)| = 1$, misalkan $K(x, y) = \{k\}$. Maka:
\[
p(y \mid x) = p(k).
\]
Namun dari keamanan sempurna, $p(y \mid x) = p(y)$, yang nilainya tidak bergantung pada $x$.
Artinya, setiap kunci $k \in \mathcal{K}$ yang memetakan $x$ ke $y$ memiliki probabilitas yang sama dengan $p(y)$.
Karena pemetaan $y \mapsto k$ adalah bijeksi untuk $x$ tertentu, maka $p(k)$ bernilai konstan untuk seluruh $k \in \mathcal{K}$.
Karena $\sum_{k \in \mathcal{K}} p(k) = 1$, maka:
\[
p(k) = \frac{1}{|\mathcal{K}|}, \quad \forall k \in \mathcal{K}.
\]
Ini membuktikan syarat (1).
\end{solusi}

\begin{soalshannon}{Soal 19 [Shannon - Sifat Entropi \& Subaditivitas]}
Misalkan $X$ dan $Y$ adalah dua variabel acak diskret pada ruang berhingga.
\begin{enumerate}[label=(\alph*)]
    \item Buktikan identitas aturan rantai entropi:
    \[
    H(X, Y) = H(Y) + H(X \mid Y).
    \]
    \item Buktikan ketaksamaan subaditivitas:
    \[
    H(X \mid Y) \le H(X) \iff H(X, Y) \le H(X) + H(Y)
    \]
    menggunakan sifat kecekungan (\textit{concavity}) fungsi $\log_2(t)$ (Ketaksamaan Jensen), serta tentukan syarat perlu dan cukup agar kesamaan berlaku.
\end{enumerate}
\end{soalshannon}

\begin{solusi}
\textbf{(a)} Berdasarkan definisi entropi bersama:
\begin{align*}
H(X, Y) &= -\sum_{x \in \mathcal{X}} \sum_{y \in \mathcal{Y}} p(x, y) \log_2 p(x, y) \\
&= -\sum_{x \in \mathcal{X}} \sum_{y \in \mathcal{Y}} p(x, y) \log_2 \big(p(y) p(x \mid y)\big) \\
&= -\sum_{x \in \mathcal{X}} \sum_{y \in \mathcal{Y}} p(x, y) \log_2 p(y) - \sum_{x \in \mathcal{X}} \sum_{y \in \mathcal{Y}} p(x, y) \log_2 p(x \mid y) \\
&= -\sum_{y \in \mathcal{Y}} \left(\sum_{x \in \mathcal{X}} p(x, y)\right) \log_2 p(y) + \sum_{y \in \mathcal{Y}} p(y) \left(-\sum_{x \in \mathcal{X}} p(x \mid y) \log_2 p(x \mid y)\right) \\
&= -\sum_{y \in \mathcal{Y}} p(y) \log_2 p(y) + \sum_{y \in \mathcal{Y}} p(y) H(X \mid Y = y) \\
&= H(Y) + H(X \mid Y).
\end{align*}

\textbf{(b)} Tinjau selisih:
\begin{align*}
H(X) - H(X \mid Y) &= -\sum_{x} p(x) \log_2 p(x) - \left(-\sum_{x, y} p(x, y) \log_2 p(x \mid y)\right) \\
&= -\sum_{x, y} p(x, y) \log_2 p(x) + \sum_{x, y} p(x, y) \log_2 \frac{p(x, y)}{p(y)} \\
&= \sum_{x, y} p(x, y) \log_2 \frac{p(x, y)}{p(x) p(y)}.
\end{align*}
Terapkan Ketaksamaan Jensen pada fungsi cekung $\log_2(t)$:
\begin{align*}
\sum_{x, y} p(x, y) \log_2 \frac{p(x) p(y)}{p(x, y)} &\le \log_2\left(\sum_{x, y} p(x, y) \frac{p(x) p(y)}{p(x, y)}\right) \\
&= \log_2\left(\sum_{x} p(x) \sum_{y} p(y)\right) \\
&= \log_2(1 \cdot 1) = 0.
\end{align*}
Akibatnya:
\[
-\sum_{x, y} p(x, y) \log_2 \frac{p(x, y)}{p(x) p(y)} \le 0 \implies \sum_{x, y} p(x, y) \log_2 \frac{p(x, y)}{p(x) p(y)} \ge 0.
\]
Maka:
\[
H(X) - H(X \mid Y) \ge 0 \iff H(X \mid Y) \le H(X).
\]
Tambahkan $H(Y)$ pada kedua ruas:
\[
H(Y) + H(X \mid Y) \le H(X) + H(Y) \iff H(X, Y) \le H(X) + H(Y).
\]
Karena fungsi $\log_2(t)$ cekung murni (\textit{strictly concave}), kesamaan berlaku jika dan hanya jika argumennya konstan untuk seluruh pasangan dengan $p(x, y) > 0$:
\[
\frac{p(x) p(y)}{p(x, y)} = 1 \iff p(x, y) = p(x) p(y), \quad \forall x, y.
\]
$\therefore$ Kesamaan berlaku jika dan hanya jika variabel acak $X$ dan $Y$ saling bebas (\textit{independen}).
\end{solusi}

\begin{soalshannon}{Soal 20 [Shannon - Enkoding Huffman]}
Diberikan alfabet sumber $X = \{a, b, c, d, e\}$ dengan distribusi probabilitas:
\[
p(a) = 0.35, \quad p(b) = 0.25, \quad p(c) = 0.20, \quad p(d) = 0.12, \quad p(e) = 0.08.
\]
\begin{enumerate}[label=(\alph*)]
    \item Konstruksikan pohon Huffman biner dan tentukan penetapan kata-kode biner (\textit{prefix-free code}) optimal untuk setiap simbol.
    \item Hitung panjang rata-rata kode:
    \[
    L(C) = \sum_{x \in X} p(x) \ell(x).
    \]
    \item Hitung entropi $H(X)$ dan verifikasi bahwa hasil enkoding memenuhi Batas Teorema Shannon Tanpa Derau:
    \[
    H(X) \le L(C) < H(X) + 1.
    \]
\end{enumerate}
\end{soalshannon}

\begin{solusi}
\textbf{(a)} Algoritma Huffman:
\begin{itemize}[leftmargin=*]
    \item \textbf{Tahap 0 (Urutan awal):} $a(0.35), b(0.25), c(0.20), d(0.12), e(0.08)$.
    \item \textbf{Tahap 1:} Gabungkan $d(0.12)$ dan $e(0.08) \to \{d, e\}(0.20)$.\\
    Node aktif: $a(0.35), b(0.25), c(0.20), \{d, e\}(0.20)$.
    \item \textbf{Tahap 2:} Gabungkan $c(0.20)$ dan $\{d, e\}(0.20) \to \{c, d, e\}(0.40)$.\\
    Node aktif: $\{c, d, e\}(0.40), a(0.35), b(0.25)$.
    \item \textbf{Tahap 3:} Gabungkan $a(0.35)$ dan $b(0.25) \to \{a, b\}(0.60)$.\\
    Node aktif: $\{a, b\}(0.60), \{c, d, e\}(0.40)$.
    \item \textbf{Tahap 4:} Gabungkan $\{a, b\}(0.60)$ dan $\{c, d, e\}(0.40) \to \text{Root}(1.00)$.
\end{itemize}
Penetapan bit biner (Cabang kiri $= 0$, Cabang kanan $= 1$):
\begin{itemize}
    \item Dari Root: Cabang kanan $\to \{a, b\}$ (prefiks $1$), Cabang kiri $\to \{c, d, e\}$ (prefiks $0$).
    \item Node $\{a, b\}$: $a \to 10$, $b \to 11$.
    \item Node $\{c, d, e\}$: $c \to 00$, $\{d, e\} \to 01$.
    \item Node $\{d, e\}$: $d \to 010$, $e \to 011$.
\end{itemize}
Tabel Enkoding Huffman Bebas-Prefiks:
\[
\begin{array}{cccc}
\toprule
\text{Simbol } x & \text{Peluang } p(x) & \text{Kata-Kode } C(x) & \text{Panjang } \ell(x) \\
\midrule
a & 0.35 & 10 & 2 \\
b & 0.25 & 11 & 2 \\
c & 0.20 & 00 & 2 \\
d & 0.12 & 010 & 3 \\
e & 0.08 & 011 & 3 \\
\bottomrule
\end{array}
\]

\textbf{(b)} Panjang rata-rata kode:
\begin{align*}
L(C) &= \sum_{x \in X} p(x) \ell(x) \\
&= 0.35(2) + 0.25(2) + 0.20(2) + 0.12(3) + 0.08(3) \\
&= 0.70 + 0.50 + 0.40 + 0.36 + 0.24 \\
&= 2.20 \text{ bit/simbol}.
\end{align*}

\textbf{(c)} Entropi sumber $H(X)$:
\begin{align*}
H(X) &= -\sum_{x \in X} p(x) \log_2 p(x) \\
&= -\big(0.35\log_2 0.35 + 0.25\log_2 0.25 + 0.20\log_2 0.20 + 0.12\log_2 0.12 + 0.08\log_2 0.08\big) \\
&\approx -\big(0.35(-1.5146) + 0.25(-2.0000) + 0.20(-2.3219) + 0.12(-3.0589) + 0.08(-3.6439)\big) \\
&\approx 0.5301 + 0.5000 + 0.4644 + 0.3671 + 0.2915 \\
&= 2.1531 \text{ bit}.
\end{align*}
Verifikasi Batas Teorema Shannon Tanpa Derau:
\[
H(X) \le L(C) < H(X) + 1 \iff 2.1531 \le 2.20 < 3.1531.
\]
Ketaksamaan tersebut terbukti benar dan terpenuhi secara presisi.
\end{solusi}

\end{document}
